\section{Historical experiments retained outside the main argument}
\label{supp:history}
The original 28-group ledger~\cite{history} remains authoritative; Table~\ref{tab:ledger} maps each group to its narrative destination. Earlier manuscripts, figures, tables and raw/capsule references remain unchanged in the repository. This supplement reorganizes their presentation without changing facts or eligibility and does not duplicate every archive. Historical per-task values remain in the pinned original evidence.

\subsection{Prediction fidelity is not reference recovery}
The B570 prospective predictor matched all 72 native arrays on a reserved 24-task/three-repeat cohort, while only 60/72 matched the ordinary reference. An initial RTX target-only predictor obtained 24/36 and failed its gate; the reserved trial did not proceed. A later independent-proposal replay obtained 72/72 reproduction, again with 60/72 reference agreement. These studies explain implementation behavior but do not establish an inexpensive online repair.

A separate RTX continuation intervention reproduced 129/129 outputs across 43 tasks and three repeats in each condition. It repaired HumanEval 15 while regressing 25, leaving aggregate reference agreement unchanged at 117/129. MBPP304 remained divergent while changing a token at index 33. The draft-free U2 bridge reproduced the intervention's complete 57/57 endpoints on 19 exposed HumanEval tasks without a draft, state transplant or future answer tape. That is an output-reproduction result, not proven full-state equivalence or a speed measurement.

The offline selector study found no supported evidence-acquisition advantage. RTX direct enumeration and screening both abstained after 114 record equivalents beyond common acquisition. B570 direct verification used 72 while screening used 144. These counts are not latency and do not constitute an implemented adaptive repair selector. The new paper does not adopt them as the central method.

\subsection{Local crosses and their limits}
The historical RTX state/path cross used 24 natural donors, 48 forks and 24 own-path shams on four probes. One selected boundary followed the current path; another displayed a donor-by-path interaction; two were local discrete nulls. Twelve prompt-KV comparisons matched, while 70/72 generated-prefix K/V tensor comparisons differed. These local findings motivated phase-aware control but did not establish which individual CUDA kernel caused every response change.

A B570 sustained alternate-path study found no changed immediate decisions across 195 paired events in 72 forks, yet 18/72 full trajectories from six tasks later diverged. Later alternate blocks continued using the alternate path. The experiment does not isolate a single changed block followed by restoration. The selected historical B570 state crosses also included recipes that changed both UB and output reservation. A donor effect under identical columns was a stronger local control, but no general stochastic or opaque-context restoration theorem followed.

\subsection{Earlier CUDA qualification and width boundaries}
The original Qwen3 expansion used 24 MBPP inputs, four arms and four repetitions: 384 requests, 96 qualified outputs per arm. The reported generation ratios were $R/C=1.170124$ and $M/C=1.169238$; resident-request ratios were $R/C=1.741569$, $R/M=1.682560$ and $M/C=1.035071$. Seven matched-request tasks regressed and one round's aggregate reversed below parity. These figures describe the original input/timer contract and are not the current main performance table.

A separate joint $(k,U)$ sweep tested (1,2), (3,4), (5,6), (7,8) at caps 128/256/512 on 11 inputs. MBPP229 failed in the two wider geometries, while matched controls passed. An outside-panel HumanEval 18 failure retained its own denominator. The maximum observed reference length was 261, so cap 512 was not evidence of a 512-token complete trajectory. A selected donor-state EOG cross at MBPP229 explained a local boundary without a universal cause of all width failures.

A fixed-UB4 study instead widened $k$ while preserving qualification. Its $k=7$ versus $k=3$ generation/request ratios were 0.829442/0.889868; versus matched target-only, 0.997404/0.950353. Draft proposals rose from 3,528 to 5,392 while accepted draft tokens rose from 2,772 to 3,192. In separate captured sessions, target callers fell from 305 to 199 but micrographs rose from 305 to 397. Those observations have different timing/instrumentation scopes and are not additive GPU cost measurements. A resident reset/order study supplied finite serial request reuse evidence, not concurrency or production throughput.

\subsection{Earlier Vulkan repair and its different endpoint}
The older Qwen3/Vulkan repair retained one live target context, sampled a pending token once and switched continuation through a phase-binding interface. Its published endpoint qualification and raw-suffix accounting are distinct from the later Llama Q6 campaigns. The four-arm formal matrix had 119/120 records after one old-frontend infrastructure interruption. R/M/C retained 15 complete task pairs, while old-N1 comparisons had 14.

The candidate's generation/request ratios were 0.282588/0.201290 versus ordinary and 0.285761/0.205704 versus matched target-only. It improved on a much slower old frontend but did not improve deployment relative to those baselines. Endpoint IDs totaled 1,150 while raw IDs totaled 1,176; all 26 suffix IDs remained charged. These results must not be replaced by the later, numerically larger but still negative Llama Q6 ratios as if they were a paired software optimization curve.

\subsection{Other platforms and control costs}
The M4 runtime study executed 60 tasks, three conditions and three repeats: 540 requests. Fixed/always speculation and the guard each passed 180 direct-ID comparisons under its own post-prefill runtime contract. Always speculation achieved about 1.346894 speed; the matched secondary guard/always comparison found the guard slower by approximately 2.71\%. A hindsight oracle restricted to two fixed actions provided only a 1.000692 factor over always speculation. This narrow result does not bound arbitrary intra-request control policies.

The 358H stateful study's CPU/NPU drafts each passed 90 direct-ID comparisons but had speeds approximately 0.73392/0.35166. Holding measured target work and committed progress fixed, a prompt-level nonnegative-cost elimination bound yielded ceilings 0.8372/0.8448, still below parity. The resulting 16.28\%/15.52\% target-work reduction requirements are conditional budgets, not achieved optimizations or a proof against NPUs. Early platform tables also mixed backend, operating system, model representation, energy and text/ID contracts. None is used as a uniform cross-hardware ranking in the reconstructed paper.
